%% RECONSTRUCTED from the submitted PDF (rgsa_submission.pdf) for latexdiff only.
%% The original submitted .tex was not retained; text and tables follow the PDF.
\documentclass[11pt]{article}

\usepackage[T1]{fontenc}
\usepackage[utf8]{inputenc}
\usepackage{lmodern}
\usepackage{microtype}
\usepackage[margin=1in]{geometry}
\usepackage{amsmath,amssymb,amsthm,mathtools}
\usepackage{booktabs}
\usepackage{array}
\usepackage{algorithm2e}
\usepackage{xcolor}
\usepackage{url}
\usepackage{hyperref}
\usepackage{placeins}
\hypersetup{colorlinks=true,linkcolor=blue,citecolor=blue,urlcolor=blue}

\newtheorem{lemma}{Lemma}
\newtheorem{remark}{Remark}
\newcommand{\R}{\mathbb{R}}
\newcommand{\Z}{\mathbb{Z}}
\newcommand{\B}{\{0,1\}}
\newcommand{\res}{\boldsymbol{\Delta}}
\newcommand{\RGSA}{\textsc{RGSA}}
\DeclareMathOperator*{\argmin}{arg\,min}

\title{Residual-Guided Simulated Annealing for\\
Multi-Dimensional Subset Sum}
\author{Sreekuttan L. S.\\
Bloq Quantum, India\\
\texttt{sreekuttan@bloq.in}}
\date{}

\begin{document}
\maketitle

\begin{abstract}
The multi-dimensional subset sum problem, also known as Market Split, asks for a binary vector satisfying several integer subset-sum constraints simultaneously. We introduce Residual-Guided Simulated Annealing (RGSA), a cardinality-preserving method whose central proposal replaces a selected item by the available vector nearest to the ideal compensator $w_i-\res$, where $\res=Ax-b$ is the current residual. We evaluate RGSA in three stages of increasing scope on the same QOBLIB instances. First, a controlled benchmark of 160 runs (eight instances, four RGSA-family configurations, five common seeds) shows that RGSA lowers the Euclidean residual $R=\|Ax-b\|_2$ relative to blind simulated annealing; on the $(m,D)=(15,200)$ flagship instance the full RGSA pipeline reaches mean residual $100.37$ against $104.06$ for RGSA alone, and no configuration certifies exact feasibility. Second, we extend the study to a common 25-second time budget and add six external exact and heuristic solvers, giving 1,554 end-to-end runs; among these, 600 complete, 416 return feasible incumbents, and 538 exceed the budget. Third, we extend the budget to 300 seconds and record intermediate checkpoints at 25, 60, 120, and 300 seconds over 1,554 independently validated runs. At 300 seconds the median residual is $38.00$ for the selected RGSA configuration, $42.25$ for blind simulated annealing, and $36.54$ for the full RGSA pipeline, while a QUBO simulated-annealing sampler reaches only $263.46$. To our knowledge this is the first controlled multi-seed heuristic scaling study on the QOBLIB Market Split classes $m\geq12$, and RGSA is a new cardinality-preserving geometric swap heuristic for this problem: it consistently produces low-residual, independently verified near-feasible assignments in seconds at the $m=15$ frontier, where exact solvers are currently impractical. These descriptive results provide a reproducible classical baseline and a candidate incumbent for exact or hybrid search; they are not proofs of infeasibility, do not certify feasibility, and do not replace exact solvers.
\end{abstract}

\noindent\textbf{Keywords:} multi-dimensional subset sum; Market Split; simulated annealing; local search; QOBLIB; heuristic optimization

\section{Introduction}\label{sec:introduction}
Market Split is the feasibility problem of selecting a subset whose aggregate attributes simultaneously attain prescribed targets. Cornu\'ejols and Dawande introduced deliberately small, hard $0$--$1$ systems to expose weaknesses in general-purpose integer programming \cite{cornuejols1998}. Subsequent work used basis reduction, lattice enumeration, branch-and-cut, and meet-in-the-middle techniques \cite{aardal1999,aardal2000,wassermann2002,wassermann2021,wassermann2026,kempkekoch2026}. The problem has renewed significance as one of the ten classes in the Quantum Optimization Benchmarking Library (QOBLIB), a common testbed intended to make comparisons between classical and quantum optimization methods reproducible \cite{koch2025}.

\paragraph{Related work: violation-guided and Lagrangian local search.}
Using constraint violations to guide move selection is established in several heuristic traditions. Guided Local Search~\cite{gls1999} penalizes features associated with local minima. Feasibility Jump~\cite{fj2023} is a particularly close conceptual precedent: it performs stochastic guided local search over a Lagrangian relaxation, selecting variable values that reduce weighted constraint violation and adapting weights for persistent violations without requiring an LP relaxation. Shift-and-Propagate~\cite{berthold2015} likewise selects shifts using violated-constraint information. Mean-field guided simulated annealing~\cite{isakov2003} uses aggregate order parameters to direct configuration search. For number partitioning, Johnson et al.~\cite{johnson1991} evaluated simulated annealing with 2-swap neighborhoods, but their proposals were not based on a multidimensional residual target. These works establish that neither guided annealing nor violation-aware local search is new in general.

\RGSA is a problem-specific specialization of this broader violation-guided local-search idea to fixed-cardinality binary partitions. For a selected item $i$, it forms the geometric ideal compensator $w^*=w_i-\res$ and chooses the available swap partner nearest to $w^*$ in the $m$-dimensional attribute space. This differs from Feasibility Jump's single-variable jump values and from penalty-based feature guidance: RGSA preserves cardinality at every move and uses the residual equation itself to define a nearest-vector proposal. To our knowledge, this particular cardinality-preserving geometric swap formulation has not been reported for Market Split/QOBLIB. The Lagrangian projection and two-item exchange used in the complete implementation are supporting components, not the central novelty claim.

The QOBLIB instances are unusual benchmarks for approximation heuristics. Their stated objective is feasibility, so an exact solver either supplies a certificate $Ax=b$ or remains unresolved. Yet a near-feasible binary vector is also scientifically useful: its distance from the target measures how effectively a heuristic navigates a dense, frustrated landscape, gives an incumbent to a hybrid exact method, and supplies a transparent classical baseline for approximate quantum methods. Classical partitioning heuristics and simulated annealing have a long history \cite{kirkpatrick1983,johnson1991,mertens1999}; evolutionary approaches have also been studied for the related multidimensional two-way number partitioning problem (MDTWNPP) \cite{kojic2010,santucci2021}. To our knowledge, however, no published study reports a controlled multi-seed heuristic scaling experiment on the QOBLIB Market Split classes $m\geq12$.

The gap is particularly visible at the exact frontier. In the QOBLIB convention an instance is considered solved once a single feasible binary vector $x$ with $Ax=b$ is found. Under this convention, recent lattice enumeration solves the QOBLIB Market Split classes up to $m=14$ on a standard computer (and enumerates all feasible solutions up to $m=9$), while $m=15$ remains the first unsolved class under the reported limits \cite{wassermann2026}; a GPU Schroeppel--Shamir implementation solves classes up to $m=11$ but remains exponential \cite{kempkekoch2026}. This paper asks a complementary question: when an exact feasible certificate is out of reach within a short budget, how close to feasibility can a lightweight classical heuristic get? RGSA does not certify feasibility and does not ``solve'' instances in the QOBLIB sense; it produces low-residual near-feasible assignments.

Our contributions are threefold.
\begin{enumerate}
  \item We provide a controlled multi-seed heuristic study of selected published QOBLIB Market Split classes $m=12$--$15$, using 5 common seeds on 8 instances and explicitly defined residual, violation, feasibility, runtime, and memory metrics.
  \item We formulate and evaluate \emph{Residual-Guided Simulated Annealing (RGSA)}, a geometric, direction-guided, cardinality-preserving swap neighborhood. Given a removed vector $w_i$ and residual $\res$, RGSA searches for the available item nearest to $w_i-\res$ and mixes this proposal with random swaps inside a Metropolis annealing schedule.
  \item We evaluate Lagrangian warm-starting and deterministic two-item polishing as supporting components. Their combination with RGSA is the full RGSA pipeline, and the ablation identifies residual guidance as the dominant empirical contributor.
  \item We report three experiments of increasing scope on the same instances: the original 160-run benchmark; a common 25-second budget comparison that adds external exact and heuristic solvers; and a 300-second budget extension with intermediate time checkpoints.
\end{enumerate}

All numerical tables and summaries are generated directly from the validated experiment artifacts.

\section{Problem Formulation and Exact-Method Barrier}\label{sec:problem}
\subsection{Market Split}
Let $A=[w_1,\ldots,w_N]\in\Z_{\geq0}^{m\times N}$ contain one attribute vector $w_i\in\Z_{\geq0}^{m}$ per item, and let $b\in\Z_{\geq0}^{m}$. The feasibility version of Market Split is
\begin{equation}
  \text{find }x\in\B^N\quad\text{such that}\quad Ax=b.
  \label{eq:msp}
\end{equation}
For a heuristic candidate we define the signed residual and energy
\begin{equation}
  \res(x)=Ax-b,\qquad R(x)=\|\res(x)\|_2,
  \qquad E(x)=\tfrac12R(x)^2.
  \label{eq:residual}
\end{equation}
Exact feasibility is equivalent to $R(x)=0$. We search within the fixed-cardinality slice
\begin{equation}
  \mathcal{X}_q=\{x\in\B^N:\mathbf{1}^{\mathsf T}x=q\},
  \label{eq:cardinality}
\end{equation}
where $q$ is supplied by the instance-generation convention or chosen near $N/2$. This is an algorithmic restriction, not part of the general definition in Eq.~\eqref{eq:msp}; all reported feasibility checks are made against $Ax=b$ itself.

Decision feasibility is NP-complete already for $m=1$ by the classical subset-sum reduction. Multiple rows couple several subset sums through a shared binary vector. Dense positive coefficients make the induced QUBO matrix $A^{\mathsf T}A$ nearly complete, while equality targets leave little tolerance for a locally plausible but globally inconsistent choice.

\subsection{QOBLIB parameterization}
The QOBLIB generated classes use $N=10(m-1)$ columns for $m\geq8$ and nonnegative integer coefficients bounded by a range parameter $D\in\{50,100,200\}$. File names have the form \texttt{ms\_mm\_DDD\_sss.dat}; the first field records $m$, the second $D$, and the final field the instance seed or identifier. The right-hand sides are based on half row sums, with coefficients adjusted around a planted near-half-cardinality assignment so that each published benchmark is feasible. Thus the largest class in this study has $m=15$ and $N=140$.

\subsection{Collapse of the continuous relaxation}
The following elementary fact explains why a natural root relaxation gives almost no information on balanced instances.

\begin{lemma}[center-point relaxation collapse]\label{lem:collapse}
Suppose $b=\tfrac12A\mathbf{1}+\varepsilon$ with $\|\varepsilon\|_\infty\le\tfrac12$ (rounding tolerance). Then $\bar{x}=\tfrac12\mathbf{1}$ is feasible for the box relaxation $x\in[0,1]^N$, attains residual $\|\varepsilon\|_p$ for the $\ell_p$-norm objective $\|Ax-b\|_p$, and satisfies $\mathbf{1}^{\mathsf T}\bar{x}=N/2$. Consequently, the relaxation has optimal value at most $\|\varepsilon\|_p$; when this is small relative to typical integer residuals, the continuous relaxation provides almost no discriminating signal.
\end{lemma}
\begin{proof}
Direct substitution gives
$A\bar{x}-b=\tfrac12A\mathbf{1}-(\tfrac12A\mathbf{1}+\varepsilon)=-\varepsilon$ and
$\mathbf{1}^{\mathsf T}\bar{x}=N/2$. The bound follows immediately.
\end{proof}

The lemma does not assert uniqueness: when $N>\operatorname{rank}(A)$ the relaxed zero-residual face can be high-dimensional. That qualification strengthens the computational diagnosis. Many fractional points achieve the same near-zero bound, even though the integer feasibility question may remain difficult. At the root node, branch-and-bound therefore sees a negligible lower bound and no objective gap that distinguishes promising binary assignments. Dense rows weaken propagation, similar columns create symmetric branches, and exact equality permits pruning only after enough variables have been fixed. Cutting planes can help, but the tree can still grow exponentially; this behavior is documented both for the original hard instances and recent QOBLIB studies \cite{cornuejols1998,vogel2012,wassermann2026}. For the QOBLIB instances, the right-hand sides satisfy $b_i=\lfloor\tfrac12\sum_j a_{ij}\rfloor$, so $\|\varepsilon\|_\infty\le\tfrac12$ and the continuous relaxation achieves residual $\|\varepsilon\|_2\approx1.9$ on the flagship instance---negligible compared with typical integer residuals above $100$.

Lattice methods address the integrality directly. Basis reduction followed by enumeration has produced the strongest CPU results and currently reaches QOBLIB $m=14$ \cite{aardal1999,wassermann2002,wassermann2021,wassermann2026}. Schroeppel--Shamir trades exponential time for reduced, but still exponential, memory and admits effective GPU acceleration \cite{kempkekoch2026}. RGSA is not a replacement for these exact methods: it sacrifices certification to obtain a useful incumbent with linear instance storage.

\section{Residual-Guided Simulated Annealing and Supporting Components}\label{sec:algorithm}
\subsection{Core method: RGSA}
RGSA is the principal method studied in this paper. It searches the fixed-cardinality slice by exchanging one selected item with one unselected item. The current residual determines a target attribute vector for the inserted item; a nearest-vector proposal exploits that target, while random proposals and Metropolis acceptance preserve exploration. The Lagrangian projection and two-item exchange described below are optional supporting components rather than prerequisites for RGSA.

For the current selected set $S$, let $\res=\sum_{k\in S}w_k-b$. If $i\in S$ is removed and $j\in\bar S$ inserted, exact cancellation requires
\begin{equation}
 \res-w_i+w_j=0
 \quad\Longleftrightarrow\quad
 w_j=w_i-\res=:w^*.
 \label{eq:ideal}
\end{equation}
RGSA therefore chooses, with probability $p_g$, the available vector nearest to $w^*$ and otherwise proposes a uniformly random insertion. This is the geometric residual-guided swap at the center of the method.

\subsection{Supporting component: Lagrangian warm start}
A nonnegative multiplier $\lambda\in\R_+^m$ projects each column to the scalar score $z_i=\lambda^{\mathsf T}w_i$. Selecting the $q$ smallest scores produces a cardinality-feasible partition. Its residual updates the multiplier, after which clipping and normalization project it back to the simplex. The method retains the best partition encountered because sorting makes the map $\lambda\mapsto x$ discontinuous.

\begin{algorithm}[t]
\caption{Lagrangian Dual Projection (Supporting Component)}\label{alg:stage1}
\KwIn{$A=[w_1,\ldots,w_N]$, $b$, cardinality $q$, iterations $L$, step $\alpha_0$, floor $\varepsilon$}
\KwOut{Best projected partition $x^{(1)}$}
$\lambda\leftarrow\mathbf{1}/m$; $R_{\min}\leftarrow+\infty$; $x^{(1)}\leftarrow\mathbf{0}$\;
\For{$t\leftarrow0$ \KwTo $L-1$}{
  \ForEach{$i\in\{1,\ldots,N\}$}{
    $z_i\leftarrow\lambda^{\mathsf T}w_i$\;
  }
  $S\leftarrow$ indices of the $q$ smallest entries of $z$\;
  $x\leftarrow\mathbf{1}_S$; $\res\leftarrow Ax-b$\;
  \If{$\|\res\|_2<R_{\min}$}{
    $R_{\min}\leftarrow\|\res\|_2$; $x^{(1)}\leftarrow x$\;
  }
  $\alpha_t\leftarrow\alpha_0/\sqrt{t+1}$\;
  $\widetilde\lambda\leftarrow\max\{\lambda+\alpha_t\res,\varepsilon\mathbf{1}\}$ componentwise\;
  $\lambda\leftarrow\widetilde\lambda/(\mathbf{1}^{\mathsf T}\widetilde\lambda)$\;
}
\Return{$x^{(1)}$}\;
\end{algorithm}

The Lagrangian warm start is best interpreted as residual-sensitive projection learning rather than a claim of solving a smooth dual problem: median selection is piecewise constant, and the best-so-far safeguard is essential.

\subsection{RGSA implementation and acceptance rule}
The implementation uses the target in Eq.~\eqref{eq:ideal}. The guided proposal is a problem-specific instantiation of informed neighborhood generation, not a claim that violation-guided annealing is new. It is related in spirit to Guided Local Search, Feasibility Jump, and guided annealing~\cite{gls1999,fj2023,isakov2003}, but differs in using the exact residual equation to form a nearest-vector swap while preserving cardinality at every step. Mixing this proposal with uniform random insertion protects exploration when $w^*$ lies far outside the observed point cloud.

\begin{algorithm}[t]
\caption{RGSA: Residual-Guided Simulated Annealing}\label{alg:stage2}
\KwIn{$A,b$, initial $x^{(1)}$, steps $K$, guide probability $p_g$, temperatures $T_0,T_f$}
\KwOut{Best annealed partition $x^{(2)}$}
$S\leftarrow\{i:x^{(1)}_i=1\}$; $\bar S\leftarrow\{1,\ldots,N\}\setminus S$\;
$\res\leftarrow\sum_{i\in S}w_i-b$; $R\leftarrow\|\res\|_2$\;
$x^{(2)}\leftarrow x^{(1)}$; $R_{\min}\leftarrow R$; $T\leftarrow T_0$\;
$\gamma\leftarrow(T_f/T_0)^{1/K}$\;
\For{$k\leftarrow1$ \KwTo $K$}{
  choose $i$ uniformly from $S$\;
  draw $u\sim\operatorname{Uniform}(0,1)$\;
  \eIf{$u<p_g$}{
    $w^*\leftarrow w_i-\res$\;
    $j\leftarrow\argmin_{h\in\bar S}\|w_h-w^*\|_2$\;
  }{
    choose $j$ uniformly from $\bar S$\;
  }
  $\res'\leftarrow\res+w_j-w_i$; $R'\leftarrow\|\res'\|_2$; $\delta\leftarrow R'-R$\;
  \If{$\delta<0$ \textbf{or} $\operatorname{Uniform}(0,1)<\exp(-\delta/\max\{T,\varepsilon\})$}{
    $S\leftarrow(S\setminus\{i\})\cup\{j\}$; $\bar S\leftarrow(\bar S\setminus\{j\})\cup\{i\}$\;
    $\res\leftarrow\res'$; $R\leftarrow R'$\;
    \If{$R<R_{\min}$}{
      $R_{\min}\leftarrow R$; $x^{(2)}\leftarrow\mathbf{1}_S$\;
    }
  }
  $T\leftarrow\gamma T$\;
}
\Return{$x^{(2)}$}\;
\end{algorithm}

The Metropolis energy in Algorithm~\ref{alg:stage2} is the residual norm rather than its square. This choice keeps temperature magnitudes interpretable and matches the implementation evaluated here. Kirkpatrick et al. established the general annealing principle \cite{kirkpatrick1983}, and Johnson et al. evaluated it for number partitioning \cite{johnson1991}.

\subsection{Supporting component: deterministic two-item exchange polishing}
Annealing is followed by deterministic descent. First, improving one-for-one exchanges are accepted. Then the best exchange of two selected for two unselected items is applied; exhaustive passes continue until no strict improvement remains or a pass limit is reached. This proves only local optimality in the examined neighborhoods, not global feasibility.

\begin{algorithm}[t]
\caption{Two-Item Exchange Polish (Supporting Component)}\label{alg:stage3}
\KwIn{$A,b$, initial $x^{(2)}$, maximum passes $P$, tolerance $\tau$}
\KwOut{Locally polished partition $x^{(3)}$}
$S\leftarrow\{i:x^{(2)}_i=1\}$; $\bar S\leftarrow\{1,\ldots,N\}\setminus S$\;
$\res\leftarrow\sum_{i\in S}w_i-b$; $R\leftarrow\|\res\|_2$\;
\ForEach{$i\in S$}{
  \ForEach{$j\in\bar S$}{
    $R'\leftarrow\|\res+w_j-w_i\|_2$\;
    \If{$R'<R-\tau$}{
      exchange $i$ and $j$ between $S$ and $\bar S$\;
      $\res\leftarrow\res+w_j-w_i$; $R\leftarrow R'$; \textbf{break}\;
    }
  }
}
$p\leftarrow0$; $\textit{improved}\leftarrow\mathrm{true}$\;
\While{$p<P$ \textbf{and} $\textit{improved}$}{
  $\textit{improved}\leftarrow\mathrm{false}$; $R_{\mathrm{pass}}\leftarrow R$\;
  $(i_1^*,i_2^*,j_1^*,j_2^*)\leftarrow\varnothing$\;
  \ForEach{$\{i_1,i_2\}\in\binom{S}{2}$}{
    \ForEach{$\{j_1,j_2\}\in\binom{\bar S}{2}$}{
      $\res'\leftarrow\res+w_{j_1}+w_{j_2}-w_{i_1}-w_{i_2}$\;
      \If{$\|\res'\|_2<R_{\mathrm{pass}}-\tau$}{
        store this exchange as $(i_1^*,i_2^*,j_1^*,j_2^*)$\;
        $R_{\mathrm{pass}}\leftarrow\|\res'\|_2$\;
      }
    }
  }
  \If{$(i_1^*,i_2^*,j_1^*,j_2^*)\neq\varnothing$}{
    apply the stored two-item exchange; recompute $\res$ and $R$\;
    $\textit{improved}\leftarrow\mathrm{true}$\;
  }
  $p\leftarrow p+1$\;
}
$x^{(3)}\leftarrow\mathbf{1}_S$; \Return{$x^{(3)}$}\;
\end{algorithm}

\subsection{Full RGSA pipeline}
The complete implementation optionally composes the supporting components around RGSA: Lagrangian projection supplies the initial partition, RGSA performs the main annealing search, and deterministic one- and two-item exchanges polish its best incumbent. We call this combination the \emph{full RGSA pipeline}. For reproducibility, its stored configuration identifier in the committed benchmark artifact and code is \texttt{full\_esa}; this label is an implementation artifact and is not the scientific name of the method.

\subsection{Complexity}
Storing $A$, the current and best assignments, index sets, $\lambda$, and residual vectors requires $O(Nm)$ memory. One Lagrangian iteration costs $O(Nm+N\log N)$ with comparison sorting. In a guided RGSA step, evaluating all candidates in $\bar S$ costs $O(Nm)$; a random step costs $O(m)$. Thus $K$ annealing steps cost $O(KNm)$ in the worst case, or $O(K[p_gNm+(1-p_g)m])$ in expectation. For RGSA, guidance is $O(Nm)$ per annealing step and $O(KNm)$ over a run. Batched or indexed nearest-neighbor search may reduce constants or asymptotic query time.

A literal exhaustive two-for-two pass costs
$O\bigl(\binom{q}{2}\binom{N-q}{2}m\bigr)=O(N^4m)$ time. Pair-sum batching reduces interpreter overhead but not the number of candidate combinations. The reported memory bound excludes optional materialization of all pair-pair residuals; the implementation stores only one side's $O(N^2m)$ pair sums and can stream candidates. For the benchmark sizes considered, RGSA dominates under a small polishing-pass limit.

\section{Experimental Setup}\label{sec:setup}
\subsection{Instances and protocol}
We evaluate one selected instance per $(m,D)$ class for $m\in\{9,10,11,12,13,14,15\}$ with $D=50$, plus the $(m{=}15,D{=}200)$ flagship, giving 8 instances in total. Each of the four stochastic configurations is run with 5 common seeds (42--46), for $8\times4\times5=160$ primary runs. Parameters are fixed: $L=3000$ Lagrangian iterations, $K=150{,}000$ SA steps, $p_g=0.65$, $T_0=600$, $T_f=10^{-3}$, and one exhaustive two-item pass for the Full RGSA pipeline. Cardinality is set to $q=\lfloor N/2\rfloor$ from the instance dimensions alone; the embedded QOBLIB reference solution is not consulted when configuring the primary runs. The separate cardinality sweep is exploratory and is not included in the 160-run primary comparison.

\subsection{Metrics}
For each run we report:
\begin{align}
 R(x)&=\|Ax-b\|_2 &&\text{(Euclidean residual)},\nonumber\\
 R_{\mathrm{rel}}(x)&=100\,\|Ax-b\|_2/\max\{\|b\|_2,1\} &&\text{(relative residual, percent)},\nonumber\\
 V_{\max}(x)&=\|Ax-b\|_\infty &&\text{(maximum row violation)},\nonumber\\
 F(x)&=\mathbf{1}[Ax=b] &&\text{(exact feasibility)},\label{eq:metrics}
\end{align}
plus wall-clock time and peak resident memory. The main tables report means and sample standard deviations over the 5 common seeds; the best-residual table reports the minimum over those seeds. Pairwise comparisons use matched instance-seed pairs. The primary RGSA-versus-blind statistic is the paired residual difference, with bootstrap confidence intervals computed using a fixed seed. Exact-feasibility rates are reported as counts and, where stated, Wilson intervals.

\subsection{Baselines and ablations}
The Full RGSA pipeline is compared with four controlled configurations under the same annealing-step budget whenever applicable. Raw configuration identifiers are retained in the benchmark records; their paper-facing names are:
\begin{enumerate}
  \item \textbf{Blind SA} (`random\_blind\_sa'): random cardinality-feasible initialization, $p_g=0$, no polish;
  \item \textbf{Lagrangian + Blind SA} (`lagrangian\_blind\_sa'): Stage 1 followed by $p_g=0$ annealing;
  \item \textbf{RGSA} (`random\_guided\_sa'): random initialization, $p_g=0.65$, no polish;
  \item \textbf{Full RGSA pipeline}: Lagrangian warm start, RGSA, and deterministic polishing (stored identifier \texttt{full\_esa}).
\end{enumerate}
This design isolates the main effect of residual guidance by comparing Blind SA and RGSA, tests the warm start by comparing RGSA with the Full RGSA pipeline, and measures polishing by comparing the RGSA incumbent with the final pipeline result.

\subsection{Extended comparison: external methods and selected configuration}
The original 160-run benchmark compares only RGSA-family configurations. To place RGSA against established solvers, we extend the study to two further stages that share the same instances but add external methods under a common wall-clock budget. All comparison methods use a single thread and fixed, pinned solver versions. The remote execution environment used to run these experiments is an implementation detail and is not part of any scientific claim.

\paragraph{Selected RGSA configuration (champion).} The extended stages evaluate one representative RGSA configuration rather than all four. We refer to it as the \emph{champion}: it is the guided simulated-annealing configuration (stored identifier \texttt{full\_guided\_sa}) chosen as the RGSA finalist in a separate, frozen selection step performed before the extended runs. The champion is an RGSA method, not a new algorithm; it is reported alongside Blind SA and the full RGSA pipeline so that the extended results remain comparable with the original benchmark.

\paragraph{External comparison methods.} The extended stages add six external solvers, each applied to the same fixed-cardinality Market Split instances:
\begin{itemize}
  \item SCIP: an exact mixed-integer solver (\texttt{pyscipopt}) minimizing the integer squared residual;
  \item CP-SAT: an exact constraint solver (Google OR-Tools) minimizing the integer squared residual with a verified 64-bit objective bound;
  \item pymoo GA: a genetic algorithm (\texttt{pymoo}) with deterministic top-$q$ cardinality repair;
  \item Neal QUBO sampler: a simulated-annealing QUBO sampler (\texttt{dwave-neal}) with an instance-derived cardinality penalty;
  \item HiGHS $L_1$ and HiGHS $L_\infty$: a linear-programming solver (\texttt{highspy}) minimizing $\ell_1$ and $\ell_\infty$ surrogate objectives rather than the squared residual.
\end{itemize}
The exact solvers (SCIP, CP-SAT) can in principle return a feasibility certificate; the heuristic and surrogate solvers return approximate assignments. All methods are compared on the same independently verified Euclidean residual $R$.

\paragraph{Experiment stages and terminology.} The three stages are: (i) the original 160-run benchmark; (ii) a \emph{common 25-second experiment}, in which every method receives the same 25-second end-to-end wall-clock budget; and (iii) a \emph{300-second experiment}, which repeats the 25-second protocol with a 300-second budget and records intermediate checkpoints at 25, 60, 120, and 300 seconds. An \emph{end-to-end} (E2E) run measures the entire solver wrapper, including instance parsing and result retrieval. A run is \emph{valid} (equivalently, \emph{available}) if it returns a vector that passes independent locked-instance verification within the budget; otherwise it is a \emph{timeout} and is retained in the denominator rather than dropped or imputed. \emph{Coverage} is the fraction of planned runs that are valid. Blind SA and the champion expose a single continuous annealing trajectory, so their checkpoint values are read directly from that trajectory; the external methods are \emph{terminal-only}, meaning a result becomes available only after the solver returns, after which that returned result is carried forward to any later checkpoint.

\subsection{Hardware and reproducibility}
Experiments were executed on an Apple Silicon (arm64) system running macOS, using Python~3.9.6 with NumPy~1.26.4 and a single solver thread per run. No GPU is used by RGSA. The array payload for the largest instance ($m{=}15$, $N{=}140$) is under $50$~KB (instance, residual, partition, and multiplier arrays); Python interpreter resident memory is not represented by this payload estimate. The instance parser, solver, seed list, raw per-run records (JSON), pinned dependency declaration, and table-generation scripts accompany the paper at the review repository listed in the Data and Code Availability section. Timings are hardware-dependent and are reported for context rather than as universal speed claims.

\section{Results}\label{sec:results}

\subsection{Evidence layers and interpretation}
We report the three experiment stages separately and never pool them into a single sample. The original benchmark contains 160 runs over eight instances, four RGSA-family configurations, and five seeds. The common 25-second experiment evaluates every method (Blind SA, the champion, the full RGSA pipeline, and the six external solvers of Section 4.4) under one 25-second end-to-end budget, yielding 1,554 end-to-end runs. The 300-second experiment repeats that protocol with a longer budget and records the four time checkpoints defined in Section 4.4.

For all new analyses, the primary quality measure is the Euclidean residual
\begin{equation}
R(x)=\|Ax-b\|_2.
\end{equation}
The exact integer squared residual $Q(x)=R(x)^2$ is retained as a secondary audit metric because it is stored and verified by the result harness. Every available vector satisfies $R^2=Q$ within numerical roundoff. We also report $V_{\max}(x)=\|Ax-b\|_\infty$, coverage, and return status. Timeout and no-result rows remain in the planned denominators and are never imputed.

\subsection{Original 160-run benchmark}
Table~\ref{tab:comprehensive} presents the original controlled comparison. Each cell reports the mean $\pm$ standard deviation of $R$ over five seeds, together with mean runtime. The generated table is retained unchanged from the validated benchmark artifact.

RGSA reduces residuals relative to Blind SA across the tested classes. At $D=50$, mean RGSA residual increases from $11.171$ at $m=9$ to $24.159$ at $m=15$, while mean runtime increases from $14.21$ to $32.30$ seconds. At the flagship $(m,D)=(15,200)$ instance, RGSA has mean residual $104.060$ and the Full RGSA pipeline has mean residual $100.367$. The Full pipeline's improvement over standalone RGSA is therefore approximately $3.7$ residual units; the supporting warm start and polishing stages are useful but are not the central empirical contribution.

Table~\ref{tab:rgsa-ablation} reports the controlled flagship ablation. Across the 40 matched instance-seed pairs in the original study, RGSA is strictly better than Blind SA in 39 pairs. No run in the original 160-run benchmark achieves exact feasibility. This is an observed outcome, not a proof that any instance is infeasible.

\subsection{Common 25-second experiment}
The common 25-second experiment retains every planned E2E run. The status totals are 600 completed, 416 incumbent, and 538 timeout, for 1,016 valid results out of 1,554 planned runs. Runs that never returned a usable vector under the budget are retained as timeouts and are not dropped or imputed.

The champion is better than Blind SA on 154 of 210 matched rows, tied on none, and worse on 56; the median paired difference $R_{\mathrm{champion}}-R_{\mathrm{Blind}}$ is $-2.849$. Against Neal QUBO, the champion is better on all 210 complete pairs, with median difference $-228.406$. Full RGSA has 180 valid rows; among matched valid rows, the champion wins 64, ties 1, and loses 115, with median difference $1.595$. These comparisons are descriptive and do not convert incomplete methods into complete cases.

The final common-budget experiment therefore demonstrates why availability must be reported beside quality. SCIP, CP-SAT, and HiGHS provide no usable E2E vector under this cap; that outcome is not equivalent to a mathematical infeasibility result.

\subsection{300-second experiment}
The 300-second experiment independently validates all 1,554 runs. Blind SA and the champion provide one verified annealing trajectory; the external methods are terminal-only, so a result is carried forward only after a valid terminal return.

At 300 seconds, the median residual is $42.249$ for Blind SA, $38.000$ for the champion, and $36.543$ for Full RGSA. Neal QUBO has median residual $263.458$, and pymoo has median residual $72.272$ over its 204 available endpoint rows. The corresponding exact-integer median $Q$ values are 1,785, 1,444, 1,335.5, 69,411, and 5,223.5, respectively; $Q$ is reported only as a secondary audit quantity.

Full RGSA has the shortest median terminal return among the RGSA-family methods at approximately 20.290 seconds, whereas Blind SA and the champion intentionally use approximately 300 seconds for their anytime trajectories. Neal QUBO returns in a median 1.536 seconds. Pymoo returns 204 incumbents and times out in six rows; SCIP, CP-SAT, and both HiGHS modes time out in every planned row. Runtime is therefore a distinct outcome from residual quality and is not a universal speed ranking.

At the 300-second endpoint, the champion wins 154, ties 0, and loses 56 against Blind SA, with median residual difference $-2.849$. Against Full RGSA, the champion wins 67, ties 1, and loses 142, with median difference $1.935$. Against Neal, the champion wins all 210 complete pairs, with median difference $-228.406$. Comparisons with SCIP and CP-SAT are availability comparisons only because no comparator residual is available.

\subsection{Comparison with published exact methods}
Table~\ref{tab:rgsa-literature} contrasts approximate RGSA runtimes with published exact-solver results for matching QOBLIB classes and coefficient ranges. Exact methods return feasibility certificates when successful; RGSA returns approximate vectors. The comparison is therefore not like-for-like in solution quality. It shows the complementary role of a fast approximate incumbent, not a standardized speedup or a claim of solver dominance.

\textbf{Important distinction:} neither a nonzero RGSA residual nor an exact solver timeout establishes infeasibility. Exact feasibility is certified only by a verified vector satisfying $Ax=b$.

\input{tables/comprehensive_table.tex}
\input{tables/ablation_table.tex}

\begin{table}[htbp]
\centering
\caption{Final common-25-second E2E outcome and residual summary. Residual summaries use valid rows only; planned and unavailable rows remain explicit.}
\footnotesize
\begin{tabular}{lrrrrrr}
\toprule
Method & Planned & Valid & Timeout & Mean $R$ & Median $R$ [IQR] & Med. $V_{\max}$ \\
\midrule
Blind SA & 210 & 210 & 0 & 49.533 & 42.249 [24.215, 66.391] & 21.5 \\
Guided SA champion & 210 & 210 & 0 & 45.547 & 38.000 [21.909, 60.339] & 22.0 \\
Full RGSA pipeline & 210 & 180 & 30 & 41.691 & 32.778 [20.031, 55.460] & 20.0 \\
Neal QUBO sampler & 210 & 210 & 0 & 321.989 & 263.458 [155.702, 427.366] & 156.5 \\
pymoo genetic algorithm & 210 & 206 & 4 & 91.680 & 81.197 [43.301, 126.032] & 45.0 \\
SCIP & 210 & 0 & 210 & -- & -- & -- \\
CP-SAT & 210 & 0 & 210 & -- & -- & -- \\
HiGHS L1/Linf & 42 each & 0 & 42 each & -- & -- & -- \\
Total & 1,554 & 1,016 & 538 & -- & -- & -- \\
\bottomrule
\end{tabular}
\end{table}

\begin{table}[htbp]
\centering
\caption{300-second checkpoint availability. Each cell is verified available rows over planned rows.}
\begin{tabular}{lrrrr}
\toprule
Method & 25 s & 60 s & 120 s & 300 s \\
\midrule
Blind SA & 210/210 & 210/210 & 210/210 & 210/210 \\
Champion & 210/210 & 210/210 & 210/210 & 210/210 \\
Full RGSA & 186/210 & 209/210 & 209/210 & 210/210 \\
SCIP & 0/210 & 0/210 & 0/210 & 0/210 \\
CP-SAT & 0/210 & 0/210 & 0/210 & 0/210 \\
pymoo GA & 0/210 & 0/210 & 0/210 & 204/210 \\
Neal QUBO & 210/210 & 210/210 & 210/210 & 210/210 \\
HiGHS L1 & 0/42 & 0/42 & 0/42 & 0/42 \\
HiGHS Linf & 0/42 & 0/42 & 0/42 & 0/42 \\
\bottomrule
\end{tabular}
\end{table}

\begin{table}[htbp]
\centering
\caption{300-second endpoint residual summary. $R$ is primary; $Q$ is the secondary verified integer audit metric.}
\footnotesize
\begin{tabular}{lrrrrrr}
\toprule
Method & Planned & Available & Mean $R$ & Median $R$ [IQR] & Median $Q$ & Med. $V_{\max}$ \\
\midrule
Blind SA & 210 & 210 & 49.533 & 42.249 [24.215, 66.391] & 1,785 & 21.5 \\
Champion & 210 & 210 & 45.547 & 38.000 [21.909, 60.339] & 1,444 & 22.0 \\
Full RGSA & 210 & 210 & 43.334 & 36.543 [22.271, 56.522] & 1,335.5 & 21.0 \\
SCIP & 210 & 0 & -- & -- & -- & -- \\
CP-SAT & 210 & 0 & -- & -- & -- & -- \\
pymoo GA & 210 & 204 & 83.725 & 72.272 [39.364, 115.668] & 5,223.5 & 39.5 \\
Neal QUBO & 210 & 210 & 321.989 & 263.458 [155.702, 427.366] & 69,411 & 156.5 \\
HiGHS L1 & 42 & 0 & -- & -- & -- & -- \\
HiGHS Linf & 42 & 0 & -- & -- & -- & -- \\
\bottomrule
\end{tabular}
\end{table}

\input{tables/literature_comparison.tex}

\FloatBarrier
\section{Discussion}\label{sec:discussion}
\subsection{Residual guidance is the central contribution}
The original controlled benchmark and the two journal-level experiments support a narrow, evidence-based claim: residual-guided exchange proposals are useful on the tested QOBLIB instances. The geometric compensator $w_i-\res$ supplies directional information that Blind SA lacks. The Full RGSA pipeline and the selected champion have complementary profiles: Full RGSA often returns earlier and reaches a lower endpoint median residual, while the champion supplies a stable complete anytime trajectory.

\subsection{Availability and quality must be reported together}
The 25-second E2E study has 1,016 valid results and 538 timeouts. The 300-second study shows how availability changes with time: Full RGSA rises from 186/210 available rows at 25 seconds to 210/210 at 300 seconds, while pymoo reaches 204/210 only at the endpoint. Terminal-only external methods do not expose a valid checkpoint vector before return. A result table that reports only residuals would hide this behavior and could make incomplete methods appear artificially competitive or artificially poor.

\subsection{Complementarity with exact algorithms}
RGSA is best viewed as a low-residual incumbent generator, diagnostic baseline, and possible front end for exact search. Lattice enumeration, branch-and-cut, meet-in-the-middle, and GPU exact methods remain preferable when certification is required. A hybrid could use an RGSA assignment, residual directions, or restricted Hamming neighborhoods to prioritize exact enumeration while retaining a verifiable certificate as the final standard.

\subsection{Limitations}
RGSA is approximate: a nonzero residual proves neither infeasibility nor hardness, and a timeout proves neither solver failure nor infeasibility. The fixed-cardinality slice is an algorithmic restriction and may exclude a feasible assignment at another cardinality. The benchmark covers selected QOBLIB instances rather than all multidimensional subset-sum distributions. Runtime depends on hardware, wrapper, solver configuration, and stopping rules. External methods with no returned vector cannot be compared on residual quality, so availability is retained as a co-primary descriptive outcome. The 300-second sensitivity analysis is descriptive and was not used to retrain, reselect the champion, or create new confirmatory hypotheses.

\subsection{Future work}
Future work should compare RGSA directly with Feasibility Jump under equal time, move, initialization, and cardinality budgets; evaluate variable-cardinality neighborhoods; study row scaling, adaptive guidance, reheating, and approximate nearest-neighbor indexes; and develop a lattice--RGSA hybrid that preserves exact certification.

\section{Conclusion}\label{sec:conclusion}
RGSA converts the residual equation into a cardinality-preserving geometric swap proposal and combines it with annealing exploration. In the original 160-run benchmark, RGSA lowers residuals across the tested scales and reaches mean residual $104.060$ at $(m,D)=(15,200)$, while the Full RGSA pipeline reaches $100.367$. The original benchmark records 0/160 exact-feasibility runs.

The common 25-second experiment retains all 1,554 E2E runs: 600 completed, 416 valid incumbents, and 538 timeouts. Among valid runs, the champion has median residual $38.000$, Blind SA $42.249$, the full RGSA pipeline $32.778$ over 180 runs, Neal QUBO $263.458$, and pymoo $81.197$ over 206 runs. The 300-second experiment independently validates all 1,554 runs and confirms the residual-first picture at the endpoint: champion $38.000$, Blind SA $42.249$, full RGSA pipeline $36.543$, Neal $263.458$, and pymoo $72.272$ over 204 runs.

These results establish RGSA as a transparent approximate baseline and a candidate front end for exact search. They do not establish infeasibility, universal dominance, or a replacement for exact algorithms. The strongest conclusion is methodological: future comparisons on difficult Market Split instances should report verified Euclidean residual, robust distribution summaries, availability, runtime, and exact-feasibility status together.

\section*{Data and Code Availability}
The source code, benchmark data, and reproduction checks are available at \url{https://github.com/Sreekuttanls/rgsa-market-split}. The committed benchmark artifact \texttt{results/paper\_benchmarks/paper\_benchmark\_results.json} contains the 160 seed-level records used for the primary tables, together with the RGSA solver, the benchmark runner, the validators, the claim checks, the table-generation code, and pinned dependency declarations. The extended common-25-second and 300-second experiments are fully scripted in the same repository and can be reproduced with the provided runners; their large per-run result sets are regenerated by these scripts rather than committed. QOBLIB instances are reused under their upstream licenses from \url{https://github.com/ZIB-AOPT/QOBLIB}. A camera-ready release will add an archival repository snapshot and a persistent DOI.

\section*{Acknowledgments}
The author thanks Nikhil Londhe, researcher at Bloq Quantum, for valuable inputs and assistance in validating and checking the code and manuscript. The author used Gemini 3.7 Flash for brainstorming, outlining, and drafting assistance during manuscript preparation. All AI-assisted material was critically reviewed, revised, and verified by the author, who takes full responsibility for the study design, implementation, analysis, accuracy of citations, and final manuscript.

\section*{Declarations}
\textbf{Funding} The author declares that no funds, grants, or other support were received during the preparation of this manuscript.

\textbf{Clinical trial number} Clinical trial number: Not applicable (this computational optimization study does not involve clinical trials or human participants).

\textbf{Ethics approval, consent to participate, and consent to publish} Ethics, Consent to Participate, and Consent to Publish declarations: not applicable.

\textbf{Competing interests} The author has no competing financial or non-financial interests that are directly or indirectly related to the work submitted for publication.

\begin{thebibliography}{99}

\bibitem{cornuejols1998}
Cornu\'ejols, G., Dawande, M. (1998).
A class of hard small 0--1 programs.
In: Bixby, R.E., Boyd, E.A., R\'ios-Mercado, R.Z. (eds),
\emph{Integer Programming and Combinatorial Optimization}, LNCS 1412, 284--293.
Springer. \url{https://doi.org/10.1007/3-540-69346-7_22}.

\bibitem{aardal1999}
Aardal, K., Bixby, R.E., Hurkens, C.A.J., Lenstra, A.K., Smeltink, J.W. (1999).
Market split and basis reduction: Towards a solution of the Cornu\'ejols--Dawande instances.
In: \emph{Integer Programming and Combinatorial Optimization}, LNCS 1610, 1--16.
Springer. \url{https://doi.org/10.1007/3-540-48777-8_1}.

\bibitem{aardal2000}
Aardal, K., Hurkens, C.A.J., Lenstra, A.K. (2000).
Solving a system of linear Diophantine equations with lower and upper bounds on the variables.
\emph{Mathematics of Operations Research}, 25(3), 427--442.

\bibitem{wassermann2002}
Wassermann, A. (2002).
Attacking the market split problem with lattice point enumeration.
\emph{Journal of Combinatorial Optimization}, 6, 5--16.
\url{https://doi.org/10.1023/A:1013355015853}.

\bibitem{wassermann2021}
Wassermann, A. (2021).
Search for combinatorial objects using lattice algorithms---revisited.
In: Flocchini, P., Moura, L. (eds), \emph{Combinatorial Algorithms}, 20--33.
Springer. \url{https://doi.org/10.1007/978-3-030-79987-8_2}.

\bibitem{wassermann2026}
Wassermann, A. (2026).
Solving the market split problem with lattice enumeration.
\emph{Mathematical Programming Computation}.
\url{https://doi.org/10.1007/s12532-026-00328-z}.

\bibitem{kempkekoch2026}
Kempke, N.C., Koch, T. (2026).
A GPU accelerated variant of Schroeppel--Shamir's algorithm for solving the market split problem.
\emph{arXiv preprint arXiv:2507.05045}.
\url{https://arxiv.org/abs/2507.05045}.

\bibitem{koch2025}
Koch, T., Bernal Neira, D.E., Chen, Y., Cortiana, G., Egger, D.J., et al. (2026).
The Quantum Optimization Benchmarking Library.
\emph{Nature Computational Science}, 6, 653--671.
\url{https://doi.org/10.1038/s43588-026-00991-1}.

\bibitem{vogel2012}
Vogel, A. (2012).
Solving the market split problem via branch-and-cut.
Technical report, Zuse Institute Berlin.

\bibitem{johnson1991}
Johnson, D.S., Aragon, C.R., McGeoch, L.A., Schevon, C. (1991).
Optimization by simulated annealing: An experimental evaluation; Part II, graph coloring and number partitioning.
\emph{Operations Research}, 39(3), 378--406.
\url{https://doi.org/10.1287/opre.39.3.378}.

\bibitem{mertens1999}
Mertens, S. (1999).
A complete anytime algorithm for balanced number partitioning.
Technical report, Otto-von-Guericke University Magdeburg.

\bibitem{kojic2010}
Koji\'c, J. (2010).
A new heuristic algorithm for the multidimensional two-way number partitioning problem.
\emph{Computers \& Operations Research}, 37, 1949--1955.

\bibitem{santucci2021}
Santucci, V., Baioletti, M., Di Bari, G. (2021).
An effective evolutionary algorithm for the multidimensional two-way number partitioning problem.
\emph{Expert Systems with Applications}, 184, 115453.

\bibitem{fj2023}
Luteberget, B., Sartor, G. (2023).
Feasibility Jump: an LP-free Lagrangian MIP heuristic.
\emph{Mathematical Programming Computation}, 15, 365--388.
\url{https://doi.org/10.1007/s12532-023-00234-8}.

\bibitem{gls1999}
Voudouris, C., Tsang, E.P.K. (1999).
Guided local search and its application to the traveling salesman problem.
\emph{European Journal of Operational Research}, 113(2), 469--499.

\bibitem{berthold2015}
Berthold, T., Hendel, G. (2015).
Shift-and-propagate.
\emph{Journal of Heuristics}, 21(1), 73--106.

\bibitem{isakov2003}
Isakov, S.V., Zinn-Justin, P. (2003).
Guided simulated annealing method for optimization problems.
\emph{arXiv preprint cond-mat/0302137}.

\bibitem{kirkpatrick1983}
Kirkpatrick, S., Gelatt, C.D., Vecchi, M.P. (1983).
Optimization by simulated annealing.
\emph{Science}, 220(4598), 671--680.
\url{https://doi.org/10.1126/science.220.4598.671}.

\end{thebibliography}

\end{document}
